% =====================================================================
%  4 — La dette : un emprunt, pas une faute
% =====================================================================
\begin{frame}{Un emprunt, pas une faute}
  \vspace{0.30cm}
  {\fontsize{11}{15}\selectfont\color{ink2}%
    La dette technique n'est pas une faute morale : c'est un
    \textbf{emprunt} que l'on contracte, parfois à raison, pour livrer plus
    vite.\par}
  \vspace{0.45cm}
  \begin{columns}[t,onlytextwidth]
    \column{0.485\textwidth}
      \card{4.05cm}{La métaphore financière}{%
        \textbf{Principal} : l'écart entre le code tel qu'il est et tel
        qu'il devrait être.\\[5pt]
        \textbf{Intérêts} : le surcoût payé à chaque intervention tant que
        le principal n'est pas remboursé.\\[5pt]
        \textbf{Remboursement} : il porte un nom, le \emph{refactoring}.}
    \column{0.485\textwidth}
      \card{4.05cm}{Un arbitrage permanent}{%
        La dette n'est presque jamais un problème purement technique :
        c'est un arbitrage entre\\[5pt]
        le \emph{time-to-market} attendu par le marché,\\[5pt]
        et la soutenabilité du logiciel : coût de maintenance, risque de
        régression, fiabilité.}
  \end{columns}
\end{frame}
